%% snd_detection.tex -- subfunctionalization, neofunctionalization and dosage
%% effects among the copies of a duplicated gene.

\begin{recipe}[label=rec:snd]{Subfunctionalization, Neofunctionalization and
  Dosage-Effect Detection (SND)}

\recipepurpose{After a duplication, decide what became of the copies: whether a
  subset of them divided the ancestral expression between them
  (subfunctionalization), whether they kept its direction but amplified it
  (dosage effect), or whether one of them expresses where the ancestor did not
  (neofunctionalization). Each is a separate test with its own criterion.}

\nomaintainer

\begin{requiredinputs}
  \item An ancestral expression vector per family. In practice this is the
        family's orthologs-only centroid --- the conserved reference, not the
        all-genes mean.
  \item The paralogs' expression vectors, as columns, and a gene-to-family
        mapping.
  \item The angle between each gene and its family's ancestor, in radians. No
        routine computes these; the filters take them as given.
\end{requiredinputs}

\begin{prerequisites}
  \dependson{sub:normalization}
  \dependson{sub:family-centroids}
  \dependson{sub:gene-family-detection}
\end{prerequisites}

\begin{workflow}
  \step{Compute each paralog's angle to its family's ancestor, and pick a
        threshold on that angle. Subfunctionalization keeps the copies
        \emph{at or above} it, the dosage effect those \emph{at or below}.}
  \step{Prefilter with \code{filter\_paralogs\_by\_pattern\_subfunctionalization}
        or \code{\ldots\_dosage\_effect}. Each returns one bitmask per family,
        naming the copies that pattern will consider at all.}
  \step{Size the search with \code{calc\_work\_arr\_paralog\_subsets\_size} and
        use the \code{max\_subset\_size} it returns --- it caps the value you
        asked for to what the filtered family can support.}
  \step{Run \code{detect\_subfunctionalization} or
        \code{detect\_dosage\_effect}. Both search subsets of the filtered
        copies: the first for a subset whose sum lands within
        \optn{rdi\_threshold} of the ancestor, the second for one that stays
        within \optn{max\_angle} of it while exceeding its magnitude by
        \optn{gain\_gamma}.}
  \step{Decode each result with \code{mask\_check\_state}, which reports
        whether a given 1-based gene index is a member of that subset.}
  \step{Test neofunctionalization separately with
        \code{detect\_neofunctionalization}, on \emph{RAP-projected unit}
        vectors. It returns one boolean per gene and axis, so it names the
        tissues that changed.}
\end{workflow}

\begin{note}
  The three tests do not share an input space. Subfunctionalization and the
  dosage effect add expression vectors up, so they need the raw vectors with
  their magnitudes intact. Neofunctionalization asks only about direction, so
  it needs vectors projected into the RAP and scaled to unit length, and
  rejects anything outside $[-1, 1]$. Mixing the two is the most common way to
  misuse this module.
\end{note}

\examplescript{python/examples/snd_detection.py}{%
  Builds one synthetic family containing each pattern exactly once and runs all
  three detectors over it, so the output can be checked against what was
  planted. The repository ships no ortholog/paralog designation, so no bundled
  family could serve.}

\begin{codefragment}[language=Python,caption={Fragments to edit: the threshold, and each pattern's criterion}]
import numpy as np
from tensor_omics import (
    filter_paralogs_by_pattern_subfunctionalization,
    calc_work_arr_paralog_subsets_size,
    detect_subfunctionalization, mask_check_state,
    mask_chunk_count)

# --- angles to the ancestor are YOURS to supply ------------
angles = np.arccos(cosines)        # radians, 0 <= a <= pi
threshold = float(np.median(angles))

n_chunks = mask_chunk_count(n_genes)
masks = filter_paralogs_by_pattern_subfunctionalization(
    angles, threshold, n_families, gene_to_fam, n_chunks)

# --- one family at a time: its ancestor, its mask ----------
mask = np.asarray(masks)[:, family - 1]
sizing = calc_work_arr_paralog_subsets_size(
    max_subset_size, n_genes, mask)     # it CAPS the size

norms = np.linalg.norm(genes, axis=0)
res = detect_subfunctionalization(
    ancestor, genes, mask, sizing["max_subset_size"],
    rdi_threshold,                      # residual tolerated
    norms, (np.argsort(norms) + 1).astype(np.int32))

subsets = np.asarray(res["work_arr_paralog_subsets"])
for k in range(res["n_results"]):       # decode the bitmasks
    picked = [i for i in range(1, n_genes + 1)
              if mask_check_state(subsets[:, k], i)]
\end{codefragment}

\begin{parameters}
  \param{threshold}{The angle at which the prefilters split wide-angle copies
    from narrow-angle ones}{A single value for \emph{all} families. The methods
    document describes a per-family median; that is not expressible here, so
    pick a global quantile of the angle distribution and say which}
  \param{rdi\_threshold}{How far a subfunctionalization subset's sum may land
    from the ancestor}{The distance below which you consider the ancestral
    expression reconstructed. Larger admits looser partitions}
  \param{max\_angle}{How far a dosage subset's sum may point away from the
    ancestor}{\textbf{Defaults to $\pi$, which is no constraint at all.} The
    pattern is defined by small angles, so set this deliberately}
  \param{gain\_gamma}{How much a dosage subset must exceed the ancestor's
    magnitude}{Default \code{0.1}, i.e. about 10\% more. Raise it to demand a
    clearer amplification}
  \param{max\_subset\_size}{Largest number of copies combined into one
    candidate}{Use the value \code{calc\_work\_arr\_paralog\_subsets\_size}
    returns. The search is combinatorial, so a cap around 15 is advisable for
    large families}
\end{parameters}

\begin{expectedresults}
  On the example's planted family of five copies, each detector finds exactly
  what was placed there: subfunctionalization returns the pair whose sum
  reproduces the ancestor with residual $0.000$; the dosage effect returns the
  pair summing to $1.50\times$ the ancestor's magnitude; and
  neofunctionalization flags the new-domain copy. It also flags both
  subfunctionalized copies --- see the interpretation, that is correct and not
  a defect.
\end{expectedresults}

\begin{pitfall}[Pitfall: giving neofunctionalization raw expression vectors]
  \code{detect\_neofunctionalization} takes RAP-projected \emph{unit} vectors,
  for the ancestors, the genes and the per-axis thresholds alike, and rejects
  values outside $[-1, 1]$ with
  \code{ToxInputError: invalid input arguments (argument 'ancestors')}. Project
  with \code{omics\_vector\_RAP\_projection} and scale with
  \code{normalize\_unit\_length} first --- which works \emph{in place} and
  returns nothing, so use the array you passed in, not the return value.
\end{pitfall}

\begin{pitfall}[Pitfall: an ancestor expressed evenly everywhere]
  A gene expressed uniformly across all axes lies along the space diagonal, so
  its RAP projection is the origin and it has no direction to normalize:
  \code{normalize\_unit\_length} then fails with
  \code{Division by zero encountered}. Neofunctionalization simply cannot be
  evaluated against such an ancestor --- there is no ancestral preference for a
  copy to depart from. Detect this case and report it rather than working
  around it.
\end{pitfall}

\begin{pitfall}[Pitfall: leaving the dosage angle at its default]
  \optn{max\_angle} defaults to $\pi$, which admits every direction, so the
  dosage search reduces to ``any subset that is large enough''. The pattern is
  defined by copies pointing the \emph{same way} as the ancestor. Set the angle
  yourself, and record what you set.
\end{pitfall}

\begin{pitfall}[Pitfall: one threshold, two inclusive bounds]
  The prefilters keep angles $\geq$ the threshold for subfunctionalization and
  $\leq$ it for the dosage effect, so a copy sitting exactly at the threshold
  passes \emph{both} --- with the median as threshold and an odd number of
  copies, the middle one always does. That is harmless, but it means the two
  candidate sets are not a partition of the family.
\end{pitfall}

\begin{interpretation}
  Read the three results together; they are not mutually exclusive, and a copy
  may register in more than one. A copy that took over part of the ancestral
  expression has by definition changed direction, so neofunctionalization flags
  it too. What separates the patterns is the evidence each requires: only
  subfunctionalization asks that the copies \emph{reconstruct} the ancestor
  between them, only the dosage effect asks that they preserve its direction
  while amplifying it, and neofunctionalization asks nothing about the other
  copies at all. A copy flagged by neofunctionalization alone --- with no
  subset containing it reconstructing the ancestor --- is the clean case.

  A subfunctionalization result is a claim about a \emph{set}, not about single
  genes: it says these copies together account for the ancestral expression.
  The residual is how well they do so, and belongs in any report of the result.

  Because the search is over subsets, absence of a result is weak evidence. It
  may mean no such division occurred, or that the prefilter removed a
  participant, or that \optn{max\_subset\_size} stopped the search below the
  size needed. Record all three settings with any negative result.
\end{interpretation}

\begin{references}
  \reference{Force, Lynch, Pickett, Amores, Yan \& Postlethwait (1999).
    Preservation of duplicate genes by complementary, degenerative mutations.
    \emph{Genetics} 151(4):1531--1545.}
  \reference{Innan \& Kondrashov (2010). The evolution of gene duplications:
    classifying and distinguishing between models. \emph{Nature Reviews
    Genetics} 11:97--108.}
\end{references}

\end{recipe}
